% Options for packages loaded elsewhere
\PassOptionsToPackage{unicode}{hyperref}
\PassOptionsToPackage{hyphens}{url}
\documentclass[
 12pt,
]{ctexart}
\usepackage{xcolor}
\usepackage{amsmath,amssymb}
\setcounter{secnumdepth}{-\maxdimen} % remove section numbering
\usepackage{iftex}
\ifPDFTeX
 \usepackage[T1]{fontenc}
 \usepackage[utf8]{inputenc}
 \usepackage{textcomp} % provide euro and other symbols
\else % if luatex or xetex
 \usepackage{unicode-math} % this also loads fontspec
 \defaultfontfeatures{Scale=MatchLowercase}
 \defaultfontfeatures[\rmfamily]{Ligatures=TeX,Scale=1}
\fi
\usepackage{lmodern}
\ifPDFTeX\else
 % xetex/luatex font selection
\fi
% Use upquote if available, for straight quotes in verbatim environments
\IfFileExists{upquote.sty}{\usepackage{upquote}}{}
\IfFileExists{microtype.sty}{% use microtype if available
 \usepackage[]{microtype}
 \UseMicrotypeSet[protrusion]{basicmath} % disable protrusion for tt fonts
}{}
\makeatletter
\@ifundefined{KOMAClassName}{% if non-KOMA class
 \IfFileExists{parskip.sty}{%
 \usepackage{parskip}
 }{% else
 \setlength{\parindent}{0pt}
 \setlength{\parskip}{6pt plus 2pt minus 1pt}}
}{% if KOMA class
 \KOMAoptions{parskip=half}}
\makeatother
\usepackage{longtable,booktabs,array}
\usepackage{calc} % for calculating minipage widths
% Correct order of tables after \paragraph or \subparagraph
\usepackage{etoolbox}
\makeatletter
\patchcmd\longtable{\par}{\if@noskipsec\mbox{}\fi\par}{}{}
\makeatother
% Allow footnotes in longtable head/foot
\IfFileExists{footnotehyper.sty}{\usepackage{footnotehyper}}{\usepackage{footnote}}
\makesavenoteenv{longtable}
\setlength{\emergencystretch}{3em} % prevent overfull lines
\providecommand{\tightlist}{%
 \setlength{\itemsep}{0pt}\setlength{\parskip}{0pt}}
\usepackage{bookmark}
\IfFileExists{xurl.sty}{\usepackage{xurl}}{} % add URL line breaks if available
\urlstyle{same}
\hypersetup{
 hidelinks,
 pdfcreator={LaTeX via pandoc}}

\author{SCX}
\date{2026}

\begin{document}

\section{Theorem 3: Unidentifiability of Noise vs.~Learnable
Difficulty}\label{theorem-3-unidentifiability-of-noise-vs.-learnable-difficulty}

\begin{quote}
\textbf{Core Claim}: onlyfrom observationsdata \((x, y, \{f_m(x)\})\)
Nonecannot distinguishLabelnoise from sample intrinsicDifficultydegree. SCX framework的Assumptions (A1)-(A6)
is NOT arbitrarytechniqueConditions, andisbreakthiskindunidentifiability的minimumSufficient Condition. 

\textbf{Revision History}: 2026-06-27 --- First version complete. 
\end{quote}

\begin{center}\rule{0.5\linewidth}{0.5pt}\end{center}

\subsection{1 Theorem Statement (Theorem
Statement)}\label{ux5b9aux7406ux9648ux8ff0-theorem-statement}

\subsubsection{1.1 Notation and Setup}\label{ux7b26ux53f7ux4e0eux8bbeux5b9a}

沿using Theorem 1 Symbol/Concept体系, 补充withlower记号: 

\begin{longtable}[]{@{}
 >{\raggedright\arraybackslash}p{(\linewidth - 2\tabcolsep) * \real{0.5000}}
 >{\raggedright\arraybackslash}p{(\linewidth - 2\tabcolsep) * \real{0.5000}}@{}}
\toprule\noalign{}
\begin{minipage}[b]{\linewidth}\raggedright
Symbol/Concept
\end{minipage} & \begin{minipage}[b]{\linewidth}\raggedright
Meaning
\end{minipage} \\
\midrule\noalign{}
\endhead
\bottomrule\noalign{}
\endlastfoot
\(\mathcal{X}\) & Input Space \\
\(\mathcal{Y}\) & Label Space, \(|\mathcal{Y}| = K\) \\
\(y^* \in \mathcal{Y}\) & trueLabel (潜atvariable, \textbf{unobservable}) \\
\(y \in \mathcal{Y}\) & observationLabel \\
\(\{f_m\}_{m=1}^M\) & \(M\) 专家model \\
\(\mathcal{D}_{\text{obs}}\) &
observationData Distribution, \((X, Y, \{f_m(X)\}) \sim \mathcal{D}_{\text{obs}}\) \\
\(\mathcal{S}\) & State Space \\
\(s: \mathcal{X} \to \mathcal{S}\) & Statusseparate配function \\
\(\eta(x) = \mathbb{P}(Y \neq Y^* \mid X = x)\) & 输入-dependence的noise率 \\
\end{longtable}

\textbf{Latent Variable Model}: Data Generationbywithlower潜atprocess决determine: 

\begin{enumerate}
\def\labelenumi{\arabic{enumi}.}
\tightlist
\item
 trueLabel \(y^*\) by \(x\) throughnot知 Oracle \(f^*\)
 generate: \(y^* = f^*(x)\)
\item
 observationLabel \(y\) cancanbecausenoisebias离 \(y^*\)
\item
 专家prediction \(\{f_m(x)\}\) is \(x\) andtrainingdata的function
\item
 Status \(s(x)\) is \(\mathcal{X}\) to \(\mathcal{S}\)
 映射, describe输入的intrinsicCharacteristics
\end{enumerate}

\textbf{Observation Space}: 研究oneonlycan访问
\((x_i, y_i, f_1(x_i), \dots, f_M(x_i))\) joint distribution, Nonemethod直接observation
\(y^*\)、\(s(x)\) ornoise指示variable. 

\subsubsection{1.2 Formal Statement}\label{ux5f62ux5f0fux5316ux9648ux8ff0}

\textbf{Theorem 3 (Noise-Difficulty Unidentifiability).} to任意
\(K \geq 2\) 类separate类Problem, 任意 \(M \geq 1\) 专家, 任意has 限State Space
\(\mathcal{S}\), existenceattwo is different的Data Generationprocess \(\mathcal{P}_1\) and
\(\mathcal{P}_2\) makes: 

\textbf{(i)} in \(\mathcal{P}_1\) lower, observationLabel \(y\)
的some些incorrectby\textbf{Labelnoise} (label noise)leads to: i.e. \(y \neq y^*\), and
\(y^*\) istrueNone误的. 

\textbf{(ii)} in \(\mathcal{P}_2\) lower, has observationLabel \(y\) equalstrueLabel
\(y^*\), butsome些sample具has \textbf{intrinsicDifficultydegree} (intrinsic
hardness): 专家systemness地犯wrong并non-becauseis Labelincorrect, andisbecauseis samplethis身的ambiguity. 

\textbf{(iii)} 边际-Conditionsdecompose恒等: two processproduces完fullmutualsame的observationjoint distribution: 

\[\mathcal{P}_1\bigl(x, y, \{f_m(x)\}_{m=1}^M\bigr) = \mathcal{P}_2\bigl(x, y, \{f_m(x)\}_{m=1}^M\bigr), \quad \forall (x, y, \{f_m\}) \in \mathcal{X} \times \mathcal{Y} \times \mathcal{Y}^M\]

\textbf{(iv)} becausethis, 任意algorithm
\(\mathcal{A}: (\mathcal{X} \times \mathcal{Y} \times \mathcal{Y}^M)^n \to \{0,1\}^n\) (will
\(n\)
 observation映射to''noise/non-noise''Assessment)Nonemethodsametimecorrect识别two 世界 in the noise: attofew一 世界 in the expectincorrect率tofewis 
\(\eta\rho/2\) (where \(\eta\) is noise率, \(\rho\)
is 歧义Status的Proportion). in particular, when \(\eta > 0\) and \(\rho > 0\)
timeshouldlower boundrigorousis currently, Prooftwo 世界Nonemethodby完美distinguish. 

\begin{quote}
\textbf{Interpretation}: Theorem 3
indicate, atnot附add额outerAssumptions的Conditionslower, ``this sample的Labeliswrong的''and''this sampletoo难了 with专家all做wrong''isobservationupper等价的Proposition. SCX
frameworkmust引入Assumptions (A1)-(A6) 来breakthiskind等价ness. 
\end{quote}

\begin{center}\rule{0.5\linewidth}{0.5pt}\end{center}

\subsection{2 Proof (Proof)}\label{ux8bc1ux660e-proof}

\subsubsection{2.1 Construction Framework}\label{ux6784ux9020ux6846ux67b6}

adopt\textbf{构造ness反例method} (constructive counterexample). 构造two 世界
\(\mathcal{P}_{\text{noise}}\) (noise世界)and
\(\mathcal{P}_{\text{hard}}\) (Difficulty世界), makes它们producesmutualsame的observationdistributionbuthas 根thisis different的because果解释. 

setup二separate类Problem \(K = 2\), Label Space \(\mathcal{Y} = \{0, 1\}\). setup
\(M \geq 1\) 专家. Definitiontwo mutual斥的Status \(s_1, s_2 \in \mathcal{S}\). 

\subsubsection{2.2 世界 A: noise驱动
(Noise-Driven)}\label{ux4e16ux754c-aux566aux58f0ux9a71ux52a8-noise-driven}

in \(\mathcal{P}_{\text{noise}}\) lower: 

\textbf{Status \(s_1\)} (``清洁易懂''Status): -
sample占比: \(\mathbb{P}(X \in s_1) = \rho\), where \(\rho \in (0, 1)\) -
trueLabel恒determine: to has \(x \in s_1\), \(y^* = 0\) - noisemechanism: Labelwithprobability
\(\eta \in (0, 1/2)\) is翻转
\[\mathbb{P}(y \neq y^* \mid x \in s_1) = \eta\] -
专家准确率: to清洁sample, has 专家independent地with \(1 - \varepsilon_1\)
probabilitycorrectprediction \(y^*\)
\[\mathbb{P}(f_m(x) = y^* \mid x \in s_1, \text{clean}) = 1 - \varepsilon_1, \quad \varepsilon_1 \in (0, 1/2)\]

\textbf{Status \(s_2\)} (``Difficulty''Status): -
sample占比: \(\mathbb{P}(X \in s_2) = 1 - \rho\) - trueLabel恒determine: to has 
\(x \in s_2\), \(y^* = 0\) -
Nonenoise: \(\mathbb{P}(y \neq y^* \mid x \in s_2) = 0\) -
专家准确率relativelylow (shouldStatusto专家来说普遍Difficulty): 
\[\mathbb{P}(f_m(x) = y^* \mid x \in s_2) = 1 - \varepsilon_2, \quad \varepsilon_1 < \varepsilon_2 \leq 1/2\]

\subsubsection{2.3 世界 B: 难degree驱动
(Difficulty-Driven)}\label{ux4e16ux754c-bux96beux5ea6ux9a71ux52a8-difficulty-driven}

in \(\mathcal{P}_{\text{hard}}\) lower: 

\textbf{Status \(s_1\)} (``含难例''Status): -
sample占比: \(\mathbb{P}(X \in s_1) = \rho\) (and世界 A mutualsame) -
\textbf{ has Labeluniformis trueLabel} (Nonenoise): \(y = y^*\) to has 样this -
buttrueLabelthis身具has ambiguity: 
\[\mathbb{P}(y^* = 0 \mid x \in s_1) = 1 - \eta, \quad \mathbb{P}(y^* = 1 \mid x \in s_1) = \eta\]
where \(\eta\) and世界 A noise率mutualsame -
专家Conditions准确率: todeterminetrueLabel, 专家的predictionlinesis such aslower: - when \(y^* = 0\)
time: \(\mathbb{P}(f_m(x) = 0 \mid x \in s_1, y^* = 0) = 1 - \varepsilon_1\)
- when \(y^* = 1\)
time: \(\mathbb{P}(f_m(x) = 0 \mid x \in s_1, y^* = 1) = 1 - \varepsilon_1\)

i.e.专家totalis倾towardatpredictionCategory \(0\) (biastowardness), andtrueLabelNoneclose. becausethis: - when
\(y^* = 0\) time, 专家准确率is \(1 - \varepsilon_1\) (表现良good) - when
\(y^* = 1\) time, 专家准确率is \(\varepsilon_1\) (systemness地犯wrong)

\textbf{Status \(s_2\)} (``Difficulty''Status): -
sample占比: \(\mathbb{P}(X \in s_2) = 1 - \rho\) (and世界 A mutualsame) -
 has Labelis trueLabel: \(y = y^*\), and
\(\mathbb{P}(y^* = 0 \mid x \in s_2) = 1\) -
专家准确率: \(\mathbb{P}(f_m(x) = y^* \mid x \in s_2) = 1 - \varepsilon_2\) (and世界
A mutualsame)

\subsubsection{2.4
World Equivalence Verification}\label{ux4e16ux754cux7b49ux4ef7ux6027ux9a8cux8bc1}

现verifytwo 世界的observationjoint distribution恒等. observationdatais 
\((x, y, f_1(x), \dots, f_M(x))\), itsjoint distributionmaydecomposeis : 

\[\mathcal{P}(x, y, f_1, \dots, f_M) = \mathcal{P}(x) \cdot \mathcal{P}(y \mid x) \cdot \prod_{m=1}^M \mathcal{P}(f_m \mid x)\]

where第二步使using Theorem 1 Assumptions (A1)-(A2) (专家attodetermine \(x\)
lowerConditionsindependent)------butNote意thisat我们\textbf{does not require}shouldAssumptions普遍成立, onlyisat构造middle使its成立withproducesmost简洁的反例. 

\textbf{Step 1: 边际distribution \(\mathcal{P}(x)\)}

two世界middle
\(\mathcal{P}(X \in s_1) = \rho\), \(\mathcal{P}(X \in s_2) = 1 - \rho\). 恒等. 

\textbf{Step 2: ConditionsLabeldistribution \(\mathcal{P}(y \mid x)\)}

in \(\mathcal{P}_{\text{noise}}\) middle (世界 A, Status \(s_1\)): 
\[\begin{aligned}
\mathcal{P}_{\text{noise}}(y = 0 \mid s_1) &= \mathbb{P}(\text{clean}) \cdot \mathbb{P}(y = y^* \mid \text{clean}) + \mathbb{P}(\text{noise}) \cdot \mathbb{P}(y = y^* \mid \text{noise}) \\
&= (1 - \eta) \cdot 1 + \eta \cdot 0 = 1 - \eta \\
\mathcal{P}_{\text{noise}}(y = 1 \mid s_1) &= (1 - \eta) \cdot 0 + \eta \cdot 1 = \eta
\end{aligned}\]

in \(\mathcal{P}_{\text{hard}}\) middle (世界 B, Status \(s_1\)): 
\[\begin{aligned}
\mathcal{P}_{\text{hard}}(y = 0 \mid s_1) &= \mathbb{P}(y^* = 0 \mid s_1) = 1 - \eta \\
\mathcal{P}_{\text{hard}}(y = 1 \mid s_1) &= \mathbb{P}(y^* = 1 \mid s_1) = \eta
\end{aligned}\]

toStatus \(s_2\), two世界uniformNonenoise/歧义: 
\[\mathcal{P}(y = 0 \mid s_2) = 1, \quad \mathcal{P}(y = 1 \mid s_2) = 0\]

becausethisConditionsLabeldistributionattwo世界middle恒等. \(\checkmark\)

\textbf{Step 3: Conditions专家distribution \(\mathcal{P}(f_m \mid x)\)}

in \(\mathcal{P}_{\text{noise}}\) middle (世界 A, Status \(s_1\)): -
专家predictiondistributionbytrueLabel \(y^*\) training/习得决determine, not-dependenceobservationLabel \(y\): 
\[\mathcal{P}_{\text{noise}}(f_m = 0 \mid s_1) = \mathbb{P}(f_m = y^* \mid s_1) = 1 - \varepsilon_1\]
\[\mathcal{P}_{\text{noise}}(f_m = 1 \mid s_1) = \varepsilon_1\]

in \(\mathcal{P}_{\text{hard}}\) middle (世界 B, Status \(s_1\)): 
\[\begin{aligned}
\mathcal{P}_{\text{hard}}(f_m = 0 \mid s_1) &= \mathbb{P}(y^* = 0 \mid s_1) \cdot \mathbb{P}(f_m = 0 \mid s_1, y^* = 0) \\
&\quad + \mathbb{P}(y^* = 1 \mid s_1) \cdot \mathbb{P}(f_m = 0 \mid s_1, y^* = 1) \\
&= (1 - \eta) \cdot (1 - \varepsilon_1) + \eta \cdot (1 - \varepsilon_1) \\
&= 1 - \varepsilon_1
\end{aligned}\]

\[\begin{aligned}
\mathcal{P}_{\text{hard}}(f_m = 1 \mid s_1) &= (1 - \eta) \cdot \varepsilon_1 + \eta \cdot \varepsilon_1 = \varepsilon_1
\end{aligned}\]

toStatus \(s_2\)
two世界mutualsame: \(\mathcal{P}(f_m = 0 \mid s_2) = 1 - \varepsilon_2\), \(\mathcal{P}(f_m = 1 \mid s_2) = \varepsilon_2\). 

becausethisConditions专家distributionattwo世界middle恒等. \(\checkmark\)

\textbf{Step 4: Conditionsindependentnessverify}

attwo世界middle, 我们构造的专家attodetermine \(x\) (or等价地, todetermine
\(s(x)\))lowerConditionsindependent. thisisbecauseis : - in世界 A
middle, 专家的trainingdataindependent (byAssumptions A1 直接to出) - in世界 B
middle, 专家的predictionerrorattodetermineStatusandtrueLabellowerindependent

becausethis \(\prod_{m=1}^M \mathcal{P}(f_m \mid x)\)
的decompose成立, andeachbecause子attwo世界between匹配. \(\checkmark\)

\textbf{Step 5: joint distribution恒等}

综combined Step 1-4 得: 

\[\mathcal{P}_{\text{noise}}\bigl(x, y, \{f_m\}\bigr) = \mathcal{P}_{\text{hard}}\bigl(x, y, \{f_m\}\bigr), \quad \forall (x, y, \{f_m\})\]

two世界的observationjoint distribution完fullconsistent. \(\square\) (QED第 (iii) part)

\subsubsection{2.5
Algorithm Unidentifiability}\label{ux7b97ux6cd5ux4e0dux53efux533aux5206ux6027}

by第 (iii) part, two世界producesobservationupperindistinguishable的joint distribution. setup \(\mathcal{A}\)
is 任意algorithm, 使using \(n\) observationsample
\(D_n = \{(x_i, y_i, \{f_m(x_i)\})\}_{i=1}^n\) 输出each sample的noiseAssessment
\(z_i \in \{0, 1\}\). 

DefinitiontwokindAssessmentMeaning: - in \(\mathcal{P}_{\text{noise}}\) lower: \(z_i = 1\)
representation''shouldsampleLabelby翻转 (noise)'' - in \(\mathcal{P}_{\text{hard}}\)
lower: \(z_i = 1\) representation''shouldsampleLabelcorrectbut专家system犯wrong (Difficulty)''

byat \(\mathcal{A}\) only-dependence \(D_n\), and \(D_n\)
attwo 世界 in the distribution完fullmutualsame: 

\[\mathcal{L}_{\mathcal{P}_{\text{noise}}}(D_n) = \mathcal{L}_{\mathcal{P}_{\text{hard}}}(D_n)\]

where \(\mathcal{L}\) representationdistribution law. becausethis \(\mathcal{A}\)
的输出distributionattwo世界middlemutualsame: 

\[\mathcal{L}_{\mathcal{P}_{\text{noise}}}(\mathcal{A}(D_n)) = \mathcal{L}_{\mathcal{P}_{\text{hard}}}(\mathcal{A}(D_n))\]

butattwo世界middle, \textbf{correct的''noise''集is different}: -
\(\mathcal{P}_{\text{noise}}\) middletruecurrentlynoisesample集
\(\mathcal{Z}_{\text{noise}}^* = \{i: y_i \neq y_i^*\} = \{i: x_i \in s_1, y_i = 1\}\)
- \(\mathcal{P}_{\text{hard}}\) middletruecurrentlynoisesample集
\(\mathcal{Z}_{\text{hard}}^* = \varnothing\) (Nonenoise)

byat \(\mathcal{A}\) Nonecannot distinguishtwo世界, its输出Marked集 \(\hat{\mathcal{Z}}\)
mustsametimeis applicable totwo世界. 令 \(a\) is algorithmat歧义subset
\(\{i: x_i \in s_1, y_i = 1\}\)
upper的expectMarked率 (should值attwo世界middlemutualsame, becauseis observationdistributionmutualsame). at世界 A
middle, this些sampleuniformis noise (correctMarkedis \(1\)), incorrect率is \((1-a)\); at世界 B
middle, this些sampleuniformnon-noise (correctMarkedis \(0\)), incorrect率is \(a\). becausethis: 

\[\max\bigl(\text{Error}_{\mathcal{P}_{\text{noise}}}(\mathcal{A}), \text{Error}_{\mathcal{P}_{\text{hard}}}(\mathcal{A})\bigr) \geq \frac{\text{Error}_{\text{noise}} + \text{Error}_{\text{hard}}}{2} \geq \frac{\eta\rho}{2}\]

where \(\eta\rho\) is歧义subsetatall样this in the 占比. shouldlower boundin \(a = 1/2\)
time达to (algorithmat歧义sampleupperis identical to随机猜测). for典型noise率 \(\eta = 0.1\)
and \(s_1\) 占比 \(\rho = 0.5\), lower boundis 
\(0.025\)------althoughsmallbutrigorousis currently, provestwo世界的not完fullmaydistinguishness. determinenessunidentifiabilityConclusion (Nonemethod完美distinguish noise fromDifficulty)still holds. \(\square\)
(QED第 (iv) part)

\subsubsection{2.6 Extension: 一般 K
类separate类}\label{ux6269ux5c55ux4e00ux822c-k-ux7c7bux5206ux7c7b}

upper述构造may直接推广to \(K > 2\) case. onlymustwillnoisemodel改is uniform匀翻转 (such as
Theorem 1 A4), 并atDifficulty世界middlewill \(y^*\)
的distributionsetup置is andnoise世界的observationLabeldistributionconsistent. Complete的推广Proofseeappendix A. 

\begin{center}\rule{0.5\linewidth}{0.5pt}\end{center}

\subsection{3 Corollary
(Corollaries): 哪些Assumptionscanbreakunidentifiability? }\label{ux63a8ux8bba-corollariesux54eaux4e9bux5047ux8bbeux80fdux6253ux7834ux4e0dux53efux8bc6ux522bux6027}

Theorem 3
indicate, notadd额outerAssumptionstimenoise andDifficultyisindistinguishable的. withlowerCorollary列出\textbf{Sufficient Condition}------each Conditionssingle独出现i.e.maybreakunidentifiability. 

\subsubsection{3.1 Corollary 1: 锚点Label (Anchor
Points)}\label{ux63a8ux8bba-1ux951aux70b9ux6807ux7b7e-anchor-points}

\textbf{Corollary 1 (Anchor Labels Break Unidentifiability).}
If tosome already知subset
\(\mathcal{X}_{\text{anchor}} \subset \mathcal{X}\), 研究onecan获取\textbf{trueLabel}
\(y^*\) (such as人工复kernel), thennoise andDifficultyin \(\mathcal{X}_{\text{anchor}}\)
uppermaydistinguish. 进一步地, If \(\mathcal{X}_{\text{anchor}}\)
覆盖了 has Status, thenglobalmaydistinguish. 

\textbf{Proof Idea}: at锚点 \(x \in \mathcal{X}_{\text{anchor}}\)
upper, comparisonobservationLabel \(y\) andtrueLabel \(y^*\): - If 
\(y \neq y^*\), thenshouldsample必however含noise - If \(y = y^*\) but专家
\(f_m(x) \neq y\), thenshouldsample必howeveris''Difficulty''的

will锚点upper的专家incorrect率outer推tonon-锚点sample, maybreakglobalunidentifiability. \(\square\)

\textbf{and SCX connection}: SCX ''cold启动协议'' (cold start
protocol)-dependenceata small amount of人工reviewsample作is 锚点来initialization化Status级noise率estimate. 

\subsubsection{3.2 Corollary 2: already知专家training域 (Known Training Domains ---
A1)}\label{ux63a8ux8bba-2ux5df2ux77e5ux4e13ux5bb6ux8badux7ec3ux57df-known-training-domains-a1}

\textbf{Corollary 2 (Known Disjoint Training Data Break
Unidentifiability).} If \(M\) 专家in \(M\)
 \textbf{disjoint}的datasubsetupperindependently trained (Assumptions
A1), and研究one知道this些training集的constitutes, thennoise andDifficultymaydistinguish. 

\textbf{Proof Idea}: disjointtraining集意味着专家的incorrect模式is\textbf{independent } (todetermine
\(x\)). atnoisesampleupper, has 专家的incorrect率sametimerises (becauseis has 专家的predictionallfortrue
\(y^*\), andobservationLabel
\(y \neq y^*\)). atDifficultysampleupper, 专家的incorrectmaymutualmutualindependent (buteach 专家的incorrect率this身morehigh). 

key的quantityDifference: 令
\(\bar{e}(x) = \frac{1}{M}\sum_m \mathbf{1}\{f_m(x) \neq y\}\)
is on average专家incorrect率. atnoisesampleupper的Conditionsdistributionis : 

\[\mathbb{E}[\bar{e}(x) \mid \text{noise}] = 1 - \frac{1 - \mathbb{E}[\bar{e}(x) \mid \text{clean}]}{K-1}\]

 (Theorem 1, Lemma
1). this relationshipatonlyhas hardness的世界middleNonemethodsametime成立. becausethis, 检验 \(\bar{e}(x)\)
的distributioniswhether满足 Lemma 1 relationship式maywithdistinguishtwo世界. \(\square\)

\textbf{importantDescription}: Corollary 2
-dependenceat研究one\textbf{确切知道}training集disjoint. If training集的relationshipnot知, thenNonemethodapplicationthis检验. 

\subsubsection{3.3 Corollary 3: noiseis independent of输入 (Noise Independent of x ---
A4)}\label{ux63a8ux8bba-3ux566aux58f0ux72ecux7acbux4e8eux8f93ux5165-noise-independent-of-x-a4}

\textbf{Corollary 3 (Input-Independent Noise Rate Breaks
Unidentifiability).} If Labelnoise率
\(\eta(x) = \mathbb{P}(Y \neq Y^* \mid X = x)\) is常数 (not-dependenceat
\(x\)), andDifficultyStatus的专家incorrect率atStatusbetweenhas Difference, thennoise andDifficultymaydistinguish. 

\textbf{Proof Idea}: If \(\eta(x) = \eta\) is global常数, thenobservationLabel \(y\)
and专家incorrect \(\{f_m(x) \neq y\}\)
'sbetween的mutualcloseness模式atnoiseandDifficulty世界middlehas this质Difference. 

atnoise世界middle, \(\bar{e}(x)\)
的expect值at清洁sampleupperlow, atnoisesampleupperhigh, butnoise事件and \(x\) independent, becausethis
\(\bar{e}(x)\) Conditionsdistributionnot随 \(x\) 变化 (边际upper), onlywhenwith \(y\)
is Conditionstimeonly thenexplicit现Difference. 

atDifficulty世界middle, \(y\) distributionmay随 \(x\)
变化, and专家incorrect率also随Status变化. through检验 \(\bar{e}(x)\) and \(y\)
iswhetheratentire Input Spaceupper具has 恒determine的joint distribution模式 (i.e.noise世界的''uniform匀noise''Characteristics), maywithdistinguishtwoone. \(\square\)

\subsubsection{3.4 Corollary 4: Statussame质ness (State Homogeneity ---
A5)}\label{ux63a8ux8bba-4ux72b6ux6001ux540cux8d28ux6027-state-homogeneity-a5}

\textbf{Corollary 4 (State Homogeneity Breaks Unidentifiability).}
If Status划separate \(\mathcal{S}\) makesateach Status \(s\)
inner, 清洁sample的expect专家incorrect率 \(\mu_s\) is常数upper bound控制 (Assumptions
A5), andStatus数 \(|\mathcal{S}| \geq 2\) 并andis differentStatusbetween \(\mu_s\)
has Difference, thennoise andDifficultymaydistinguish. 

\textbf{Proof Idea}: atTheorem 3 反例构造middle, two 世界的Status \(s_1\) and
\(s_2\) inStatussame质nessAssumptionslowerwillproduces矛盾 \(\mu_s\)
estimate. specifically, atwithlowertwokind解释middletomany一kindcan满足Statussame质nessAssumptions: 

\begin{itemize}
\item
 \textbf{noise解释} (世界 A): Status \(s_1\) 清洁incorrect率is 
 \(\varepsilon_1\), butobservationto
 \(\bar{e}(s_1) = (1-\eta)\varepsilon_1 + \eta(1-\varepsilon_1)\)------thisisnoiseand清洁的mixedcombined. such as果
 \(\eta\) isnot知的, NonemethodatthisStatuslowerdistinguishnoiseandDifficultyseparate量. 
\item
 \textbf{Difficulty解释} (世界 B): Status \(s_1\) middle \(\eta\)
 Proportionsample的trueLabelis Category 1, 专家atshould子群体upper准确率onlyis 
 \(\varepsilon_1\). thisis equivalent to认is Status \(s_1\) containstwo 子Status (Category 0 
 \(s_{1a}\) andCategory 1 \(s_{1b}\)), and它们的专家incorrect率is different. 
\end{itemize}

Statussame质nessAssumptions (A5)requireeach Statusinner清洁incorrect率uniform匀. such as果 \(s_1\)
的innerpart专家准确率atis differentCategory的samplebetweenexplicit著is different (世界 B middleis 
\(1-\varepsilon_1\) vs \(\varepsilon_1\)), then (A5)
atthisStatusupperby违反. becausethis, such as果 (A5) 成立, 世界 B 构造does not hold. \(\square\)

\subsubsection{3.5 Corollary 5: balanceerrordistribution (Balanced Error Distribution ---
A6)}\label{ux63a8ux8bba-5ux5e73ux8861ux8befux5deeux5206ux5e03-balanced-error-distribution-a6}

\textbf{Corollary 5 (Balanced Error Distribution Breaks
Unidentifiability).} If 专家的incorrectat has incorrectCategoryupperuniform匀distribution (Assumptions
A6, \(C_{\text{bal}} = 1\)), thennoise andDifficultymaydistinguish. 

\textbf{Proof Idea}: atTheorem 3 世界 B 构造middle, 专家predictionbiastowardCategory 0. when
\(y^* = 1\) time, 专家with \(1 - \varepsilon_1\) highprobabilityprediction
0, this意味着atStatus \(s_1\) middle, \(f_m\) incorrect几乎all集middleatCategory 1 →
predictionis 0 directionupper. this违anti-A6 balancerequire------专家的incorrectextremelydegreenotbalance. 

becausethis, such as果already知专家的incorrectisbalance (\(C_{\text{bal}} \approx 1\)), then世界
B 构造byexcludes. \(\square\)

\subsubsection{3.6 Corollary 6: manyStatus交叠 (Multi-State
Overlap)}\label{ux63a8ux8bba-6ux591aux72b6ux6001ux4ea4ux53e0-multi-state-overlap}

\textbf{Corollary 6 (Differential Expert Accuracy Across States Breaks
Unidentifiability).} If existenceattofewtwo Status \(s_a, s_b \in \mathcal{S}\)
makes专家in \(s_a\) upper的准确率\textbf{above}in \(s_b\)
upper的准确率, 并and研究onecanmay靠地willsampleseparate配toStatus, thennoise andDifficultymaydistinguish. 

\textbf{Proof Idea}: Theorem 3 世界 A in构造upper-dependenceatStatus \(s_1\)
sametimecontains清洁andnoisesample的mixedcombined, butitsobservationLabeldistributionandStatus \(s_2\)
is different. throughcomparisonis differentStatus'sbetween的consistency scoredistribution, maywith检验''mixedcombined''Assumptionsiswhether成立. 

specifically, setup \(\mu_s = \mathbb{E}[C \mid \text{clean}, X \in s]\). at
Theorem 1 Lemma 1 middle, noisesample的consistency scoreexpectis 
\(1 - \mu_s/(K-1)\). If insome Statusmiddleobservationtoconsistency scoremean
\(\bar{C}_s > 1 - \hat{\mu}_s/(K-1)\) (where \(\hat{\mu}_s\)
isFrom lowseparate位数estimate的清洁baseline), thenshouldStatusmiddlemaydoes not existnoise (i.e.世界 B
类型). \(\square\)

\begin{center}\rule{0.5\linewidth}{0.5pt}\end{center}

\subsection{4 and SCX Assumptions体系 (A1-A6)
的connection}\label{ux4e0e-scx-ux5047ux8bbeux4f53ux7cfb-a1-a6-ux7684ux8054ux7cfb}

Theorem 3 及itsCorollaryreveal SCX Assumptions (A1)-(A6)
的根this作用: 它们is NOT 任意附add的techniqueConditions, andis\textbf{breakunidentifiability的minimumSufficient Condition}. 

\subsubsection{4.1
Assumptionsseparate类: 哪些Assumptionsatbreak什么}\label{ux5047ux8bbeux5206ux7c7bux54eaux4e9bux5047ux8bbeux5728ux6253ux7834ux4ec0ux4e48}

\begin{longtable}[]{@{}
 >{\raggedright\arraybackslash}p{(\linewidth - 6\tabcolsep) * \real{0.1875}}
 >{\raggedright\arraybackslash}p{(\linewidth - 6\tabcolsep) * \real{0.1875}}
 >{\raggedright\arraybackslash}p{(\linewidth - 6\tabcolsep) * \real{0.3438}}
 >{\raggedright\arraybackslash}p{(\linewidth - 6\tabcolsep) * \real{0.2812}}@{}}
\toprule\noalign{}
\begin{minipage}[b]{\linewidth}\raggedright
Assumptions
\end{minipage} & \begin{minipage}[b]{\linewidth}\raggedright
Content
\end{minipage} & \begin{minipage}[b]{\linewidth}\raggedright
break的歧义
\end{minipage} & \begin{minipage}[b]{\linewidth}\raggedright
toshouldCorollary
\end{minipage} \\
\midrule\noalign{}
\endhead
\bottomrule\noalign{}
\endlastfoot
\textbf{A1} & disjointtraining集 &
专家independentness: ensurenoisesampleupper has 专家sametime出wrong的模式anddry净sampleis different & Corollary
2 \\
\textbf{A2} & 清洁Conditionsindependent & todetermine \(x\) lower专家incorrect的independentnessdecompose & Corollary 2
的techniqueFoundation \\
\textbf{A3} & has 界loss & techniquenessAssumptions (Hoeffding
inequality must), not直接breakunidentifiability & --- \\
\textbf{A4} & uniform匀independentnoise & noise率的输入Noneclosenessanduniform匀distribution: excludes \(x\)
-dependence的noise率leads to的伪''Difficulty''模式 & Corollary 3 \\
\textbf{A5} & Statussame质ness &
Statusinner清洁incorrect率uniform匀: excludes了''intrinsicDifficultysample聚集atsame一Status''的mayness &
Corollary 4 \\
\textbf{A6} & balanceerrordistribution &
专家incorrectat has incorrectCategoryupperuniform匀: excludes了专家biastowardnessleads to的伪''Difficulty''模式 &
Corollary 5 \\
\end{longtable}

\subsubsection{4.2
minimumbreaking ness集combined}\label{ux6700ux5c0fux7834ux574fux6027ux96c6ux5408}

A1-A6 middle哪些isbreakunidentifiability的necessaryAssumptions? Analysissuch aslower: 

\textbf{Theorem 3 anti-例的抵抗ness}: Theorem 3 构造use了withlowerConditions: 1. \(K = 2\)
separate类 (most简singlecase) 2. 专家Conditionsindependent (A1/2) 3. noise率-dependenceatStatus (\(\eta\)
vs 0)------thisand A4 ''uniform匀independentnoise''has finemicro-差别 4.
Statusbetween专家incorrect率is different (\(\varepsilon_1 \neq \varepsilon_2\))

mustexcludesTheorem 3 构造, \textbf{tofewrequirewithlowerAssumptions's一}: - A1
的\textbf{disjointtraining域知识} (Corollary
2): If 我们知道专家training域确实disjoint, 世界 B 构造may违��专家independentness -
A4 \textbf{uniform匀noise} (Corollary 3): If noise率global恒determine, 世界 A 构造require
\(\eta\) in \(s_1\) and \(s_2\) middlemutualsame, 迫使heavynewdesign - A5
的\textbf{Statussame质ness} (Corollary 4): If Statusinnerincorrect率uniform匀, 世界 B \(s_1\)
inner \(y^*\) -dependence的incorrect率违反thisAssumptions - A6 \textbf{balanceerrordistribution} (Corollary
5): If 专家incorrectuniform匀distribution, 世界 B 专家biastowardnessbyexcludes

\textbf{becausethis, SCX breakunidentifiability must的minimumAssumptions集combinedis}: 
\[A_{\min} = \{A1, A4, A5\} \quad \text{or} \quad \{A1, A4, A6\} \quad \text{or} \quad \{A5, A6\} \text{ addupperStatus数 } |\mathcal{S}| \geq 2\]

i.e.tofewrequire\textbf{training域independentness} + \textbf{noiseuniform匀ness} +
\textbf{一kindincorrectuniform匀nessConditions} (StatusinnerorCategorybetween), orone\textbf{Statussame质ness} +
\textbf{incorrectbalanceness} + \textbf{manyStatusstructure}. 

\subsubsection{4.3
Assumptions违反time的退化}\label{ux5047ux8bbeux8fddux53cdux65f6ux7684ux9000ux5316}

理解each Assumptions的Roleback, maywithestablish违反eachAssumptionstime的退化模式: 

\begin{longtable}[]{@{}
 >{\raggedright\arraybackslash}p{(\linewidth - 4\tabcolsep) * \real{0.3929}}
 >{\raggedright\arraybackslash}p{(\linewidth - 4\tabcolsep) * \real{0.2857}}
 >{\raggedright\arraybackslash}p{(\linewidth - 4\tabcolsep) * \real{0.3214}}@{}}
\toprule\noalign{}
\begin{minipage}[b]{\linewidth}\raggedright
违反的Assumptions
\end{minipage} & \begin{minipage}[b]{\linewidth}\raggedright
退化is 
\end{minipage} & \begin{minipage}[b]{\linewidth}\raggedright
实际表现
\end{minipage} \\
\midrule\noalign{}
\endhead
\bottomrule\noalign{}
\endlastfoot
A1 (training域heavy叠) & 专家mutualclosenessleads to Theorem 1 exponential界消失 & SCX noise detection
F1 falls, butstillmayis superior to随机 \\
A4 (noise-dependence \(x\)) & noise率andCharacteristics难degreecoupled & SCX
的noiseestimatehas bias, butrelative排序may保留 \\
A5 (Statusdifferent质) & Statusinner清洁/noisemixedcombined & Theorem 1
not适用, requiremorefine的Status划separate \\
A6 (incorrectnotuniform衡) & noise detectionbiastoward常seeincorrectCategory & false positive集middleat特determineCategory, F1
fallsbutmaystillmayusing \\
\textbf{all违反} & 完full退化 & \textbf{Theorem 3
的unidentifiability生效, 任何algorithmNonemethoddistinguish noise fromDifficulty} \\
\end{longtable}

\subsubsection{4.4 SCX
frameworksuch as何implicit式usethis些Conditions}\label{scx-ux6846ux67b6ux5982ux4f55ux9690ux5f0fux4f7fux7528ux8fd9ux4e9bux6761ux4ef6}

at实际 SCX applicationmiddle, A1-A6
并non-totalisbright确检验的. frameworkthroughwithlowermechanism\textbf{implicit式}use了breakunidentifiability的structure: 

\begin{enumerate}
\def\labelenumi{\arabic{enumi}.}
\item
 \textbf{cold启动协议}: require人工reviewa small amount of锚点samplegenerateinitializationStatusLabel (is equivalent toCorollary
 1 锚点). 
\item
 \textbf{Statusdiscover}: throughto \(\phi(X)\) 聚类来划separateStatus, implicit含地Assumptions
 A5------聚类试图ateach 簇innerminimum化变different, 使Statussame质nessapproximate成立. 
\item
 \textbf{multi-expert聚combined}: usedisjointtraining的专家 (A1), thisrequireexperimentdesign阶段bright确separate配trainingdata. 
\item
 \textbf{统计检验}: check \(\hat{C}(s)\) 经验distributioniswhether满足 Lemma 1
 Mean Separationrelationship, If does not satisfythenMarkedis ''weakCharacteristics失效'' (Theorem 2). 
\end{enumerate}

\begin{center}\rule{0.5\linewidth}{0.5pt}\end{center}

\subsection{5
and现has 文献的connection}\label{ux4e0eux73b0ux6709ux6587ux732eux7684ux8054ux7cfb}

\subsubsection{5.1 measureerrormodel (Measurement Error
Models)}\label{ux6d4bux91cfux8befux5deeux6a21ux578b-measurement-error-models}

统计学的measureerror文献middlehas 一 经典Conclusion: at经典measureerrormodel (classical
measurement error model)lower, If nohas \textbf{verifydata} (validation
data)or\textbf{工具variable} (instrumental
variables), measureerror的distributionandtrue值的distributionindistinguishable (Fuller, 1987; Carroll
et al., 2006). 

form化地, tomodel \(Y = Y^* + U\), where \(U\) is measureerror: If onlyobservationto
\(Y\) and \(Y^*\) and \(U\) uniformnot知, thenhas None穷manyto \((Y^*, U)\)
的distribution满足observationto \(Y\)
distribution. thisis\textbf{measureerrorunidentifiability}的Standard result. 

Theorem 3
willthisresultFrom continuousmeasureerrorExtensiontoseparate类Labelnoise and专家Difficulty的to称Problemmiddle. key推广is: notonlyhas 
\(Y\) and \(Y^*\) relationship, still引入了专家prediction \(\{f_m(X)\}\)
作is 额outerobservation. 我们Proofi.e.使has this些额outerinformation, unidentifiabilitystill holds------unlessto
\(\{f_m\}\) structure施addAssumptions. 

\subsubsection{5.2
Labelnoisetheory的unidentifiability}\label{ux6807ux7b7eux566aux58f0ux7406ux8bbaux7684ux4e0dux53efux8bc6ux522bux6027}

Labelnoise文献 in the 经典resultis\textbf{noise转移矩阵unidentifiability} (Menon et al.,
2015; Patrini et al., 2017): 

to \(K\) 类separate类, noise转移矩阵 \(T \in [0,1]^{K \times K}\) 满足
\(T_{ij} = \mathbb{P}(\tilde{Y} = j \mid Y^* = i)\). atonlyobservationto
\((\tilde{Y}, X)\) Conditionslower, \(T\) andtrueConditionsdistribution
\(\mathbb{P}(Y^* \mid X)\) 联combinedisunidentifiable的------unlessaddupperSufficient Condition (such as
\(T\) toangle占优ness、existenceat锚点、orexistenceatmany noise率already知的markNoteone). 

Theorem 3
isthisresulttoward\textbf{multi-expertscenario}的推广. atmulti-expertscenariomiddle, 我们notonlyobservationto
\(\tilde{Y}\), stillobservationto
\(\{f_m(X)\}\), butunidentifiabilitystill holds. key洞察is: 专家this身的incorrect模式mayfrom
\(Y^*\) Difficulty区域, alsomayfrom \(Y^*\)
by翻转back的伪模式------thistwooneatobservationupper等价. 

\subsubsection{5.3 Dawid-Skene
的unidentifiability}\label{dawid-skene-ux7684ux4e0dux53efux8bc6ux522bux6027}

Dawid-Skene model (1979)middle, each markNoteone \(m\) has 一 confusion matrix
\(\pi_m\), trueLabel \(Y^*\)
is潜atvariable. Standard resultis: atto称ness约束 (such asfixed 一 markNoteone的confusion matrix)or锚点markNoteone (somemarkNoteone完fullmay靠)的Conditionslower, modelonly thenmay识别. 

Theorem 3 and Dawid-Skene unidentifiability的relationship: - Dawid-Skene
的unidentifiabilitycloseNote\textbf{is differentmarkNoteone的confusion matrix'sbetween的置换to称ness} (Label翻转Problem)
- Theorem 3 closeNote\textbf{noise andDifficulty'sbetween的语义indistinguishability}------thisis SCX
特has 一阶Problem - twoonecurrently交: i.e.使解决 Dawid-Skene
的Label置换Problem (through锚点markNoteone), Theorem 3 unidentifiabilitystillhoweverexistenceat

换言's, Dawid-Skene 解决的is''哪位markNoteonemay靠''的Problem, and Theorem 3
reveal的is''sampleis 什么notmay学''的Problem------backoneisfrontone的front提 (at解释专家incorrect模式'sfront, must知incorrect来源). 

\subsubsection{5.4
because果infer的视angle}\label{ux56e0ux679cux63a8ux65adux7684ux89c6ux89d2}

From because果infer的angledegree, Theorem 3
mayby理解is 一kind\textbf{because果structureunidentifiability} (causal structure
unidentifiability; Pearl, 2009): 

considertwo because果图: 

\textbf{图 A (noisebecause果图)}: 

\begin{verbatim}
S → Y* → Y ← Noise
↓
X → {f_m}
\end{verbatim}

where箭头 \(Y^* \to Y\) representationtrueLabel决determineobservationLabel, but \(Y\)
alsomaybyouterpartnoise直接impact. 

\textbf{图 B (Difficultybecause果图)}: 

\begin{verbatim}
S → Y* (= Y)
↓
X → {f_m}
\end{verbatim}

whereDifficultyStatus \(S\) sametimeimpacttrueLabeldistributionand专家准确率. 

two because果图generatemutualsame的observationdistribution
\((X, Y, \{f_m\})\), but蕴含is different的because果解释. onlyfrom observationsdataNonecannot distinguishthistwo 图------thisis典型的\textbf{马尔may夫等价类} (Markov
equivalence class)Problem. 

\subsubsection{5.5
mixedcombinedmodelandhas 限mixedcombined的may识别ness}\label{ux6df7ux5408ux6a21ux578bux4e0eux6709ux9650ux6df7ux5408ux7684ux53efux8bc6ux522bux6027}

athas 限mixedcombinedmodel (finite mixture
models)文献middle, may识别ness通常requireseparate量distribution具has 特determine的parameterform (such asexponential family), orrequireseparate量数has upper bound (Teicher,
1963; McLachlan \& Peel, 2000). 

Theorem 3 反例mayby视is 一 \textbf{twoseparate量mixedcombinedmodel}的unidentifiabilityexample: -
世界 A: separate量 1 (清洁, 占 \(1-\eta\))andseparate量 2 (noise, 占 \(\eta\)) -
世界 B: separate量 1 (\(y^*=0\), 占 \(1-\eta\))andseparate量 2 (\(y^*=1\), 占
\(\eta\))

two mixedcombinedmodelproducesmutualsame的observationjoint distribution, butseparate量has is different的语义解释. thisindicateatLabelnoisescenariomiddle, i.e.使separate量数already知 (\(K=2\)), Standard mixedcombinedmodelmay识别nessConditions (such asparameter化formalready知)stillhoweveris insufficient------stillrequiretoseparate量'sbetween的''专家prediction模式''施add额outer的structure约束. 

\subsubsection{5.6
minimum值andmaximum值: Assumptions的necessaryness谱系}\label{ux6700ux5c0fux503cux4e0eux6700ux5927ux503cux5047ux8bbeux7684ux5fc5ux8981ux6027ux8c31ux7cfb}

综combinedwithupper讨论, may得withlower\textbf{Assumptionsnecessaryness谱系}: 

\begin{verbatim}
mostweakAssumptions (unidentifiable)
├── NoneAssumptions → Theorem 3 成立: 完fullunidentifiable
├── + 锚点Label (Corollary 1) → in锚点uppermay识别
├── + disjointtraining域 (A1, Corollary 2) → throughincorrectindependentness检验may识别
├── + uniform匀noise (A4, Corollary 3) → throughincorrect率的输入not变nessmay识别
├── + Statussame质ness (A5, Corollary 4) → throughStatusinnerincorrect率uniform匀nessmay识别
├── + incorrectbalanceness (A6, Corollary 5) → through专家incorrectCategorydistributionmay识别
└── A1 + A4 + A5 (or A1 + A4 + A6) → 完fullmay识别 (SCX Standard Assumptions)
moststrongAssumptions (完fullmay识别)
\end{verbatim}

\begin{center}\rule{0.5\linewidth}{0.5pt}\end{center}

\subsection{6 实际meaning}\label{ux5b9eux9645ux610fux4e49}

\subsubsection{6.1 to SCX
framework的currentlywhenness化}\label{ux5bf9-scx-ux6846ux67b6ux7684ux6b63ux5f53ux6027ux5316}

Theorem 3 is SCX Assumptions体系provide了根thisnesscurrentlywhenness: 

\begin{enumerate}
\def\labelenumi{\arabic{enumi}.}
\item
 \textbf{A1-A6
 is NOT 过degreeAssumptions}: 它们isbreaknotmayavoid的unidentifiability must的minimum附addstructure. 任何声称canatNoneAssumptionsConditionslowerdistinguish noise fromDifficulty的Methodology, must么implicit式contains了等价Assumptions, must么existenceattheorydefect. 
\item
 \textbf{Theorems 1-3 commonsameconstitutes逻辑闭环}: 

 \begin{itemize}
 \tightlist
 \item
 Theorem 1: in A1-A6 lower, noisemaydetection (Positive Statement)
 \item
 Theorem 2: whenCharacteristicsweakensdetectioncan力 (Boundary Statement)
 \item
 Theorem 3: NoneAssumptionstimenoise andDifficultyindistinguishable (Necessity Statement)
 \end{itemize}
\item
 \textbf{cold启动协议的theoryFoundation}: SCX
 must求initialization阶段人工reviewa small amount ofsample来estimateStatus级noise率. Theorem 3
 Proofthisis\textbf{theory必must}andnon-工程妥协------nothas this 锚点, systemNonemethodbreakthroughunidentifiability. 
\end{enumerate}

\subsubsection{6.2
to竞争Methodology的批判nessAnalysis}\label{ux5bf9ux7adeux4e89ux65b9ux6cd5ux7684ux6279ux5224ux6027ux5206ux6790}

Theorem 3 provide了评价its他noise detectionMethodology的framework: 

\begin{itemize}
\item
 \textbf{loss-based Methodology} (such as MentorNet, Decoupling,
 Co-teaching): implicit式Assumptions''highloss = noise'', i.e.in Theorem 3
 middle选择了''世界 A''解释. whendata实际遵循''世界
 B''模式time, this些MethodologywillwillDifficultysample误markis noise. 
\item
 \textbf{置信learningMethodology} (Confidence Learning; Northcutt et al.,
 2021): use计数矩阵estimatenoise转移矩阵, implicit式Assumptions专家 (separate类器)的计数andnoise率'sbetweenexistenceatmaydecomposerelationship. at
 ``世界 B'' middle, shoulddecompose失效. 
\item
 \textbf{many视图Methodology} (Multi-view
 methods): -dependencemany independent视图的consistentness, is equivalent to Theorem 1 A1-A2
 Assumptions. Theorem 3 Proof, nothas this些Assumptionstimemany视图consistentnessalsoNonemethoddistinguish noise fromDifficulty. 
\end{itemize}

\subsubsection{6.3
experimentdesignRecommendation}\label{ux5b9eux9a8cux8bbeux8ba1ux5efaux8bae}

based on Theorem 3, 提出withlowerexperimentdesign原then: 

\begin{enumerate}
\def\labelenumi{\arabic{enumi}.}
\item
 \textbf{bright确Assumptions}: at论文orexperimentreportmiddlebright确列出 做Assumptions (such as''我们Assumptionsnoise and
 x independent''). Theorem 3 意味着not列出的Assumptionsthenisnot满足的识别Conditions. 
\item
 \textbf{sensitivenessAnalysis}: 检验ConclusiontoAssumptions违反的sensitiveness. 例such as, If Statussame质ness稍weak (A5
 notrigorous成立), noise detection F1 fallsmanyfew? 
\item
 \textbf{Assumptions检验}: design统计检验来verifykeyAssumptions. 例such as, using \(\chi^2\)
 检验专家incorrectiswhetheruniform匀distribution (检验 A6), orusing KS
 检验Statusbetween的consistency scoredistributionDifference (检验 A5). 
\item
 \textbf{manyMethodologycross-validation}: sametime运lines SCX and loss-based
 Methodology, willtwooneMarkedis inconsistent的sample提交人工review------thiscurrentlyis SCX
 cold启动协议的实际操作方式. 
\end{enumerate}

\begin{center}\rule{0.5\linewidth}{0.5pt}\end{center}

\subsection{appendix A: K 类推广}\label{ux9644ux5f55-ak-ux7c7bux63a8ux5e7f}

willTheorem 3 推广to \(K > 2\). adoptand二separate类is different的构造strategy------at世界 B
middle, 专家by构造is \textbf{完full随机} (notrelies on real Label), thismakesjoint distributionat任意
\(K\) lowermaintains恒等. 

\textbf{世界 A (noise)}: to has 
\(x \in s_1\), \(y^* = 0\). noisemechanism: withprobability \(\eta\) willLabeluniform匀翻转to
\(\mathcal{Y} \setminus \{0\}\): 
\[\mathbb{P}(y = c \mid \text{noise}, x \in s_1) = \frac{1}{K-1}, \quad c \neq 0\]

专家准确率: \(\mathbb{P}(f_m(x) = y^* \mid x \in s_1) = 1 - \varepsilon_1\), incorrectuniform匀distribution: 
\[\mathbb{P}(f_m = c \mid s_1) = \frac{\varepsilon_1}{K-1}, \quad c \neq 0\]

observationdistributionis : 
\[\mathcal{P}_{\text{noise}}(y = 0 \mid s_1) = 1 - \eta, \quad \mathcal{P}_{\text{noise}}(y = c \mid s_1) = \frac{\eta}{K-1}, \; c \neq 0\]
\[\mathcal{P}_{\text{noise}}(f_m = 0 \mid s_1) = 1 - \varepsilon_1, \quad \mathcal{P}_{\text{noise}}(f_m = c \mid s_1) = \frac{\varepsilon_1}{K-1}, \; c \neq 0\]

andbyat (A4) noise and专家independent, \(y\) and \(f_m\) intodetermine \(s_1\) lowerindependent. 

\textbf{世界 B (Difficulty)}: to
\(x \in s_1\), has Labeluniformis trueLabel (Nonenoise): 
\[\mathbb{P}(y^* = 0 \mid s_1) = 1 - \eta, \quad \mathbb{P}(y^* = c \mid s_1) = \frac{\eta}{K-1}, \; c \neq 0\]

\textbf{key构造}: at世界 B
middle, 专家\textbf{notlearning}------itsprediction完full随机andis independent oftrueLabel: 
\[\mathbb{P}(f_m = 0 \mid s_1) = 1 - \varepsilon_1, \quad \mathbb{P}(f_m = c \mid s_1) = \frac{\varepsilon_1}{K-1}, \; c \neq 0\]
and \(f_m \perp y^* \mid s_1\) (专家predictionandtrueLabelindependent). 

atthis构造lower: - \(y\) 边际distribution: two世界mutualsame (世界 A noiseLabeldistribution =
世界 B trueLabeldistribution) - \(f_m\) 边际distribution: two世界mutualsame (uniformis 
\(1-\varepsilon_1\) inCategory 0, \(\varepsilon_1/(K-1)\) inits他Category) -
\((y, f_m)\) joint distribution: two世界uniformis independent乘积 (世界 A middle \(y\) and \(f_m\)
because (A4) independent; 世界 B middle \(f_m\) ispure随机thereforealsoand \(y\) independent)

becausethisto任意
\(K \geq 2\), \(\mathcal{P}_{\text{noise}}(x, y, \{f_m\}) = \mathcal{P}_{\text{hard}}(x, y, \{f_m\})\)
恒成立. \(\square\)

\textbf{解释}: 世界 B
的构造middle, ``Difficulty''的语义isextreme的------shouldStatusmiddle\textbf{ has }sampleto has 专家allsame样Difficulty (专家is identical to随机猜测), None论trueLabelis什么. thistoshould了实际middle''Characteristics完fullNoneinformation''的extremely限case (Theorem
2 middle
\(\delta = 0\)). althoughhoweverextreme, but它作is \textbf{existenceatness反例}alreadysufficient: existenceattofew一 with''Difficulty''is 解释的世界, itsobservationdistributionandwith''noise''is 解释的世界完fullmutualsame, becausethistwooneNonemethodFrom observationdatamiddledistinguish. 

\begin{quote}
\textbf{FIXED (2026-06-28, DEFECT-12):} to \(K>2\)
的构造use\textbf{完full随机专家} (\(f_m \perp y^* \mid s_1\))------thisis''难degree''的一kindextreme/退化form. for
\(K=2\), 构造use现实form的难degree (专家biastowardCategory 0), but \(K>2\)
的构造require完full随机专家来ensurejoint distributionat has Categoryupper匹配. this意味着: althoughhowever\textbf{原thenupper}existenceatsomekind''难degree''distributionandnoiseindistinguishable, butcan实现thiskindindistinguishability的''难degree''Category\textbf{non-常狭narrow}. 现实世界 in the Difficultysample通常具has \textbf{structure化}的专家incorrect模式 (such as专家consistentincorrect地predictionsome 特determineincorrectCategory), thisat附add温andAssumptions (such as
A5 Statussame质nessor A6 balanceerrordistribution)lowermayandnoisedistinguishopen来. \textbf{\(K=2\)
的构造具has moststrong的practicemutualcloseness}, becauseis 它use现实form的专家biastowardness; \(K>2\)
的构造primary作is existenceatnessProof, itsextremeAssumptions (完full随机专家)at现实
MLIP/Medical Imaging/LLM scenariomiddle几乎is impossible出现. 
\end{quote}

\begin{quote}
\textbf{2026-06-27 correction}: 原版appendix A 构造middle, 世界 B
的专家attodeterminetrueLabellower具has non-平凡的准确率 \(1-\varepsilon_1\), thisleads tofor
\(K>2\) time专家边际distributionand世界 A
not匹配. correctionback的构造adopt''完full随机专家'', guarantee了to任意 \(K \geq 2\)
joint distribution的恒等ness. 
\end{quote}

\begin{center}\rule{0.5\linewidth}{0.5pt}\end{center}

\subsection{appendix B: Symbol/Concept表}\label{ux9644ux5f55-bux7b26ux53f7ux8868}

\begin{longtable}[]{@{}lll@{}}
\toprule\noalign{}
Symbol/Concept & Meaning & 首次出现 \\
\midrule\noalign{}
\endhead
\bottomrule\noalign{}
\endlastfoot
\(\mathcal{P}_{\text{noise}}\) & noise世界的Data Generationprocess & 2.2 \\
\(\mathcal{P}_{\text{hard}}\) & Difficulty世界的Data Generationprocess & 2.3 \\
\(\eta\) & noise率 (世界 A)/ Label歧义率 (世界 B) & 2.2 \\
\(\varepsilon_1\) & Status \(s_1\) upper的专家incorrect率 & 2.2 \\
\(\varepsilon_2\) & Status \(s_2\) upper的专家incorrect率 & 2.2 \\
\(\rho\) & Status \(s_1\) 边际probability & 2.2 \\
\(\mathcal{D}_{\text{obs}}\) & observationData Distribution & 1.1 \\
\(\mathcal{Z}^*\) & truenoisesample集 & 2.5 \\
\end{longtable}

\begin{center}\rule{0.5\linewidth}{0.5pt}\end{center}

\subsection{参考文献}\label{ux53c2ux8003ux6587ux732e}

\begin{enumerate}
\def\labelenumi{\arabic{enumi}.}
\item
 Fuller, W. A. (1987). \emph{Measurement Error Models}. Wiley.
\item
 Carroll, R. J., Ruppert, D., Stefanski, L. A., \& Crainiceanu, C. M.
 (2006). \emph{Measurement Error in Nonlinear Models: A Modern
 Perspective} (2nd ed.). Chapman and Hall/CRC.
\item
 Menon, A. K., Van Rooyen, B., Ong, C. S., \& Williamson, R. C. (2015).
 Learning from corrupted binary labels via class-probability
 estimation. \emph{ICML}.
\item
 Patrini, G., Rozza, A., Menon, A. K., Nock, R., \& Qu, L. (2017).
 Making deep neural networks robust to label noise: A loss correction
 approach. \emph{CVPR}.
\item
 Dawid, A. P., \& Skene, A. M. (1979). Maximum likelihood estimation of
 observer error-rates using the EM algorithm. \emph{Journal of the
 Royal Statistical Society: Series C (Applied Statistics)}, 28(1),
 20-28.
\item
 Northcutt, C. G., Jiang, L., \& Chuang, I. L. (2021). Confident
 learning: Estimating uncertainty in dataset labels. \emph{Journal of
 Artificial Intelligence Research}, 70, 1373-1411.
\item
 Pearl, J. (2009). \emph{Causality: Models, Reasoning, and Inference}
 (2nd ed.). Cambridge University Press.
\item
 Teicher, H. (1963). Identifiability of finite mixtures. \emph{The
 Annals of Mathematical Statistics}, 34(4), 1265-1269.
\item
 McLachlan, G. J., \& Peel, D. (2000). \emph{Finite Mixture Models}.
 Wiley.
\item
 SCX Theorem 1: Multi-Expert Consistency Guarantees for Label Noise
 Detection. \texttt{./01\_noise\_detection\_guarantee.md}.
\item
 SCX Theorem 2: Weak Feature Failure Lower Bound.
 \texttt{./02\_weak\_feature\_failure.md}.
\item
 SCX Data Four-Classification.
 \texttt{../../knowledge/02\_/data四separate类.md}.
\item
 SCX Data Poisoning Defense.
 \texttt{../../knowledge/02\_/data防middle毒.md}.
\end{enumerate}

\begin{center}\rule{0.5\linewidth}{0.5pt}\end{center}

\textbf{correctionDescription (2026-06-27)}: First version complete. Theorem 3 orientationis SCX
Assumptions体系的''necessaryness支柱''------Proof A1-A6
A1-A6 are not arbitrary assumptions, andisbreakunidentifiability must的minimumstructureConditions. 

\end{document}
